% SPDX-License-Identifier: CC-BY-NC-ND-4.0
% !TEX root = main.tex

\section{モナド}

\begin{definition} \label{def: Endofunctor}
$\mathscr{C}$を圏とする.\ 共変関手$F \colon \mathscr{C} \to \mathscr{C}$を\term{じこかんしゅ}{自己関手}{endofunctor}と呼ぶ.
\end{definition}

\begin{definition} \label{def: EndofunctorCategory}
$\mathscr{C}$を圏とする.\ $\End(\mathscr{C}) \coloneq \mathscr{C}^{\mathscr{C}}$としこれを\term{じこかんしゅけん}{自己関手圏}{cateogory of endofunctors}という.
\end{definition}

\begin{definition} \label{def: Monad}
$\mathscr{C}$を圏とし$T \in \End(\mathscr{C})$とする.\ 自然変換$\eta \colon \id_{\mathscr{C}} \Rightarrow T$と$\mu \colon T^2 \Rightarrow T$が次の2つの図式を可換にするとする.
% https://q.uiver.app/#q=WzAsOCxbMCwwLCJUXjMiXSxbMSwwLCJUXjIiXSxbMCwxLCJUXjIiXSxbMSwxLCJUIl0sWzMsMCwiVCJdLFs0LDAsIlReMiJdLFs0LDEsIlQiXSxbMywxLCJUXjIiXSxbMCwxLCJUXFxtdSIsMCx7ImxldmVsIjoyfV0sWzAsMiwiXFxtdSBUIiwyLHsibGV2ZWwiOjJ9XSxbMiwzLCJcXG11IiwyLHsibGV2ZWwiOjJ9XSxbMSwzLCJcXG11IiwwLHsibGV2ZWwiOjJ9XSxbNSw2LCJcXG11IiwwLHsibGV2ZWwiOjJ9XSxbNyw2LCJcXG11IiwyLHsibGV2ZWwiOjJ9XSxbNCw1LCJcXGV0YSBUIiwwLHsibGV2ZWwiOjJ9XSxbNCw3LCJUIFxcZXRhIiwyLHsibGV2ZWwiOjJ9XSxbNCw2LCJcXG1hdGhybXtpZH1fVCIsMSx7ImxldmVsIjoyfV1d
\[\begin{tikzcd}
	{T^3} & {T^2} && T & {T^2} \\
	{T^2} & T && {T^2} & T
	\arrow["{T\mu}", Rightarrow, nfold, from=1-1, to=1-2]
	\arrow["{\mu T}"', Rightarrow, nfold, from=1-1, to=2-1]
	\arrow["\mu", Rightarrow, nfold, from=1-2, to=2-2]
	\arrow["{\eta T}", Rightarrow, nfold, from=1-4, to=1-5]
	\arrow["{T \eta}"', Rightarrow, nfold, from=1-4, to=2-4]
	\arrow["{\mathrm{id}_T}"{description}, Rightarrow, nfold, from=1-4, to=2-5]
	\arrow["\mu", Rightarrow, nfold, from=1-5, to=2-5]
	\arrow["\mu"', Rightarrow, nfold, from=2-1, to=2-2]
	\arrow["\mu"', Rightarrow, nfold, from=2-4, to=2-5]
\end{tikzcd}\]
このとき組$(T, \eta, \mu)$を$\mathscr{C}$上の\term{もなど}{モナド}{monad}といい$\eta$を\term{たんいげん}{単位元}{unit},\ $\mu$を\term{じょうほう}{乗法}{multiplication}と呼ぶ.
\end{definition}

\begin{proposition} \label{prop: monad_of_adjunction}
$\mathscr{C}, \mathscr{D}$を圏とし$F \in \mathscr{D}^{\mathscr{C}},\ U \in \mathscr{C}^{\mathscr{D}}$とする.\ 随伴$F \dashv U$が存在するとき$(\eta, \varepsilon) \coloneq F \dashv U$とすると$(UF, \eta, U \varepsilon F)$は$\mathscr{C}$上のモナドになる.
\end{proposition}
\begin{proof}
$\varepsilon F \circ F \eta = \id_F$より$U \varepsilon F \circ UF \eta = \id_{UF}$が成り立ち$U \varepsilon \circ \eta U = \id_U$より$U \varepsilon F \circ \eta UF = \id_{UF}$が成り立つので次の図式が可換になる.
% https://q.uiver.app/#q=WzAsNCxbMCwwLCJVRiJdLFsxLDEsIlVGIl0sWzEsMCwiVUZVRiJdLFswLDEsIlVGVUYiXSxbMCwxLCJcXG1hdGhybXtpZH1fe1VGfSIsMSx7ImxldmVsIjoyfV0sWzAsMiwiVUYgXFxldGEiLDAseyJsZXZlbCI6Mn1dLFsyLDEsIlUgXFx2YXJlcHNpbG9uIEYiLDAseyJsZXZlbCI6Mn1dLFswLDMsIlxcZXRhIFVGIiwyLHsibGV2ZWwiOjJ9XSxbMywxLCJVIFxcdmFyZXBzaWxvbiBGIiwyLHsibGV2ZWwiOjJ9XV0=
\[\begin{tikzcd}
	UF & UFUF \\
	UFUF & UF
	\arrow["{UF \eta}", Rightarrow, nfold, from=1-1, to=1-2]
	\arrow["{\eta UF}"', Rightarrow, nfold, from=1-1, to=2-1]
	\arrow["{\mathrm{id}_{UF}}"{description}, Rightarrow, nfold, from=1-1, to=2-2]
	\arrow["{U \varepsilon F}", Rightarrow, nfold, from=1-2, to=2-2]
	\arrow["{U \varepsilon F}"', Rightarrow, nfold, from=2-1, to=2-2]
\end{tikzcd}\]

$((U \varepsilon F) \circ (U \varepsilon FUF))_X = (U \varepsilon F)_X \circ (U \varepsilon FUF)_X = U(\varepsilon_{F(X)}) \circ U(\varepsilon_{FUF(X)}) = U(\varepsilon_{F(X)} \circ \varepsilon_{FUF(X)})$である.\ また$((U \varepsilon F) \circ (UFU \varepsilon F))_X = (U \varepsilon F)_X \circ (UFU \varepsilon F)_X = U(\varepsilon_{F(X)}) \circ UFU(\varepsilon_{F(X)}) = U(\varepsilon_{F(X)} \circ FU \varepsilon_{F(X)})$である.\ $\varepsilon$は自然変換なので任意の$f \in \Mor_{\mathscr{D}}(FUF(X), F(X))$に対して次の図式が可換になる.
% https://q.uiver.app/#q=WzAsNCxbMSwwLCJGVUYoWCkiXSxbMSwxLCJGKFgpIl0sWzAsMCwiRlUoRlVGKFgpKSJdLFswLDEsIkZVKEYoWCkpIl0sWzAsMSwiZiJdLFsyLDMsIkZVKGYpIiwyXSxbMiwwLCJcXHZhcmVwc2lsb25fe0ZVRihYKX0iXSxbMywxLCJcXHZhcmVwc2lsb25fe0YoWCl9Il1d
\[\begin{tikzcd}
	{FU(FUF(X))} & {FUF(X)} \\
	{FU(F(X))} & {F(X)}
	\arrow["{\varepsilon_{FUF(X)}}", from=1-1, to=1-2]
	\arrow["{FU(f)}"', from=1-1, to=2-1]
	\arrow["f", from=1-2, to=2-2]
	\arrow["{\varepsilon_{F(X)}}", from=2-1, to=2-2]
\end{tikzcd}\]
$\varepsilon_{F(X)} \in \Mor_{\mathscr{D}}(FUF(X), F(X))$なので$\varepsilon_{F(X)} \circ \varepsilon_{FUF(X)} = \varepsilon_{F(X)} \circ FU(\varepsilon_{F(X)})$が成り立ち$U(\varepsilon_{F(X)} \circ \varepsilon_{FUF(X)}) = U(\varepsilon_{F(X)} \circ FU \varepsilon_{F(X)})$である.\ よって次の図式が可換になる.
% https://q.uiver.app/#q=WzAsNCxbMCwwLCJVRlVGVUYiXSxbMSwwLCJVRlVGIl0sWzEsMSwiVUYiXSxbMCwxLCJVRlVGIl0sWzAsMSwiVSBcXHZhcmVwc2lsb24gRlVGIiwwLHsibGV2ZWwiOjJ9XSxbMSwyLCJVIFxcdmFyZXBzaWxvbiBGIiwwLHsibGV2ZWwiOjJ9XSxbMywyLCJVIFxcdmFyZXBzaWxvbiBGIiwwLHsibGV2ZWwiOjJ9XSxbMCwzLCJVRlUgXFx2YXJlcHNpbG9uIEYiLDIseyJsZXZlbCI6Mn1dXQ==
\[\begin{tikzcd}
	UFUFUF & UFUF \\
	UFUF & UF
	\arrow["{U \varepsilon FUF}", Rightarrow, nfold, from=1-1, to=1-2]
	\arrow["{UFU \varepsilon F}"', Rightarrow, nfold, from=1-1, to=2-1]
	\arrow["{U \varepsilon F}", Rightarrow, nfold, from=1-2, to=2-2]
	\arrow["{U \varepsilon F}", Rightarrow, nfold, from=2-1, to=2-2]
\end{tikzcd}\]

したがって$(UF, \eta, U \varepsilon F)$は$\mathscr{C}$上のモナドである.
\end{proof}

\begin{definition} \label{def: monadOfAdjunction}
\cref{prop: monad_of_adjunction}におけるモナド$(UF, \eta, U \varepsilon F)$を随伴$F \dashv U$によって生成されるモナドという.
\end{definition}

\begin{definition} \label{def: EilenbergMooreCategory}
$\mathscr{C}$を圏とし$(T, \eta, \mu)$を$\mathscr{C}$上のモナドとする.\ 圏$\mathscr{C}^T$を次のようにして定める.
\begin{itemize}
	\item 対象を$A \in \mathscr{C},\ a \in \Mor_{\mathscr{C}}(T(A), A)$の組$(A, a)$で次の図式を可換にするものとする.
% https://q.uiver.app/#q=WzAsNyxbMCwwLCJBIl0sWzEsMCwiVChBKSJdLFsxLDEsIkEiXSxbMywwLCJUXjIoQSkiXSxbNCwwLCJUKEEpIl0sWzQsMSwiQSJdLFszLDEsIlQoQSkiXSxbMCwxLCJcXGV0YV9BIl0sWzEsMiwiYSJdLFswLDIsIlxcbWF0aHJte2lkfV9BIiwyXSxbMyw0LCJcXG11X0EiXSxbNCw1LCJhIl0sWzMsNiwiVChhKSIsMl0sWzYsNSwiYSIsMl1d
\[\begin{tikzcd}
	A & {T(A)} && {T^2(A)} & {T(A)} \\
	& A && {T(A)} & A
	\arrow["{\eta_A}", from=1-1, to=1-2]
	\arrow["{\mathrm{id}_A}"', from=1-1, to=2-2]
	\arrow["a", from=1-2, to=2-2]
	\arrow["{\mu_A}", from=1-4, to=1-5]
	\arrow["{T(a)}"', from=1-4, to=2-4]
	\arrow["a", from=1-5, to=2-5]
	\arrow["a"', from=2-4, to=2-5]
\end{tikzcd}\]
	\item 射$f \colon (A, a) \to (B, b)$は$\mathscr{C}$の射$f \colon A \to B$で次の図式を可換にするものとする.
% https://q.uiver.app/#q=WzAsNCxbMCwwLCJUKEEpIl0sWzEsMCwiVChCKSJdLFswLDEsIkEiXSxbMSwxLCJCIl0sWzAsMSwiVChmKSJdLFsyLDMsImYiXSxbMCwyLCJhIiwyXSxbMSwzLCJiIl1d
\[\begin{tikzcd}
	{T(A)} & {T(B)} \\
	A & B
	\arrow["{T(f)}", from=1-1, to=1-2]
	\arrow["a"', from=1-1, to=2-1]
	\arrow["b", from=1-2, to=2-2]
	\arrow["f", from=2-1, to=2-2]
\end{tikzcd}\]
\end{itemize}
この圏を$T$の\term{あいれんべるぐむーあけん}{アイレンベルグムーア圏}{Eilenberg-Moore cateogry}という.
\end{definition}

\begin{definition} \label{def: Algebra}
$\mathscr{C}$を圏とし$(T, \eta, \mu)$を$\mathscr{C}$上のモナドとする.\ アイレンベルグムーア圏$\mathscr{C}^T$の対象を\term[T-algebra]{てぃーだいすう}{$T$-代数}{$T$-algebra}と呼ぶ.
\end{definition}

\begin{definition} \label{def: forgetfulFunctor}
$\mathscr{C}$を圏とし$(T, \eta, \mu)$を$\mathscr{C}$上のモナドとする.\ 共変関手$U^T \colon \mathscr{C}^T \to \mathscr{C}$を次のようにして定める.
\begin{itemize}
	\item $(A, a) \in \mathscr{C}^T$に対して$U^T(A, a) \coloneq A$とする.
	\item $f \in \Mor_{\mathscr{C}^T}((A,a), (B,b))$に対して$U^T(f) \coloneq f$とする.
\end{itemize}
この関手を\term{ぼうきゃくかんしゅ}{忘却関手}{forgetful functor}と呼ぶ.
\end{definition}

\begin{definition} \label{def: freeFunctor}
$\mathscr{C}$を圏とし$(T, \eta, \mu)$を$\mathscr{C}$上のモナドとする.\ 共変関手$F^T \colon \mathscr{C} \to \mathscr{C}^T$を次のようにして定める.
\begin{itemize}
	\item $X \in \mathscr{C}$に対して$F^T(X) \coloneq (T(X), \mu_X)$とする.
	\item $f \in \Mor({\mathscr{C}})$に対して$F^T(f) \coloneq T(f) \in \Mor_{\mathscr{C}^T}((T(\dom(f)), \mu_{\dom(f)}), (T(\cod(f)), \mu_{\cod(f)}))$とする.
\end{itemize}
この関手を\term{じゆうかんしゅ}{自由関手}{free functor}と呼ぶ.
\end{definition}

\begin{definition} \label{def: freeAlgebra}
$\mathscr{C}$を圏とし$(T, \eta, \mu)$を$\mathscr{C}$上のモナドとする.\ $F^T$を自由関手としたとき$X \in \mathscr{C}$に対して$F^T(X)$を$X$によって定まる\term[free T-algebra]{じゆうてぃーだいすう}{自由$T$-代数}{free $T$-algebra}という.
\end{definition}

\begin{proposition} \label{prop: free_adjoint_forgetful}
$\mathscr{C}$を圏とし$(T, \eta, \mu)$を$\mathscr{C}$上のモナドとすると随伴$F^T \dashv U^T$が存在する.
\end{proposition}
\begin{proof}
任意の$X \in \mathscr{C}$に対して$U^T F^T(X) = U^T(T(X), \mu_X) = T(X)$であり任意の$f \in \Mor({\mathscr{C}})$に対して$U^T F^T(f) = U^T(T(f)) = T(f)$なので$U^T F^T = T$である.

自然変換$\varepsilon \colon F^T U^T \Rightarrow \id_{\mathscr{C}^T}$を$\varepsilon_{(A,a)} \coloneq a \in \Mor_{\mathscr{C}^T}((T(A), \mu_A), (A, a))$で定める.\ 任意の$X \in \mathscr{C}$に対して$(\varepsilon F^T \circ F^T \eta)_X = (\varepsilon F^T)_X \circ (F^T \eta)_X = \varepsilon_{F^T(X)} \circ F^T(\eta_X) = \varepsilon_{(T(X), \mu_X)} \circ T(\eta_X) = \mu_X \circ T(\eta_X) = (\mu \circ T \eta)_X = (\id_T)_X = (\id_{F^T})_X$なので$\varepsilon F^T \circ F^T \eta = \id_{F^T}$である.

任意の$(X, x) \in \mathscr{C}^T$に対して$(U^T \varepsilon \circ \eta U^T)_{(X, x)} = (U^T \varepsilon)(X, x) \circ (\eta U^T)_{(X,x)} = U^T(x) \circ \eta_X = x \circ \eta_X = \id_X = (\id_{U^T})_{(X,x)}$なので$U^T \varepsilon \circ \eta U^T = \id_{U^T}$である.

したがって$\eta$を単位,\ $\varepsilon$を余単位として$F^T \dashv U^T$である.
\end{proof}

\begin{proposition} \label{prop: monad_of_free_adjunction}
\cref{prop: free_adjoint_forgetful}における随伴$F^T \dashv U^T$によって生成されるモナドは$(T, \eta, \mu)$である.
\end{proposition}
\begin{proof}
\cref{prop: free_adjoint_forgetful}の仮定の上で$U^T F^T = T$であり任意の$X \in \mathscr{C}$に対して$(U^T \varepsilon F^T)_X = U^T(\varepsilon_{(T(X), \mu_X)}) = U^T(\mu_X) = \mu_X$であるから\cref{prop: monad_of_adjunction}により$F^T \dashv U^T$によって生成されるモナドは$(T, \eta, \mu)$である.
\end{proof}

% 随伴が与えられればモナドが構成できる. モナドが与えられればそのモナドを構成する随伴が構成できる.

\begin{definition} \label{def: KleisliCategory}
$\mathscr{C}$を圏とし$(T, \eta, \mu)$を$\mathscr{C}$上のモナドとする.\ 圏$\mathscr{C}_T$を次のようにして定める.
\begin{itemize}
	\item $\Ob(\mathscr{C}_T) \coloneq \Ob(\mathscr{C})$とする.
	\item $f \in \Mor_{\mathscr{C}_T}(A, B) \coloneq \Mor_{\mathscr{C}}(A, T(B))$とする.\ 射の合成と単位射は次のようにして定める.
	\begin{itemize}
		\item $f \in \Mor_{\mathscr{C}_T}(A, B),\ g \in \Mor_{\mathscr{C}_T}(B, C)$に対して$g \circ_{\mathscr{C}_T} f = \mu_C \circ_{\mathscr{C}} T(g) \circ_{\mathscr{C}} f$とする.
		\item $X \in \Ob_{\mathscr{C}_T}$に対して$\id_X \coloneq \eta_X$とする.
	\end{itemize}
\end{itemize}
この圏を\term{くらいすりけん}{クライスリ圏}{Kleisli category}と呼ぶ.
\end{definition}

\begin{proposition} \label{prop: kleisli_adjunction}
$\mathscr{C}$を圏とし$(T, \eta, \mu)$を$\mathscr{C}$上のモナドとする.\ 共変関手$F_T \colon \mathscr{C} \to \mathscr{C}_T$を$A \in \Ob(\mathscr{C})$に対して$F_T(A) = A$とし$f \in \Mor(\mathscr{C})$に対して$F_T(f) \coloneq \eta_{\cod(f)} \circ f$で定める.\ 共変関手$U_T \colon \mathscr{C}_T \to \mathscr{C}$を$A \in \Ob(\mathscr{C})$に対して$U_T(A) \coloneq T(A)$とし$g \in \Mor_{\mathscr{C}}(A, T(B))$に対して$U_T(g) \coloneq \mu_B \circ T(g)$とする.\ このとき$F_T \dashv U_T$が存在する.
\end{proposition}

\begin{proof}
任意の$A \in \mathscr{C}$に対して$U_T F_T(A) = U_T (A) = T(A)$であり任意の$f \in \Mor(\mathscr{C})$に対して$U_T F_T(f) = U_T(\eta_{\cod(f)} \circ f) = \mu_{\cod(f)} \circ T(\eta_{\cod(f)} \circ f) = \mu_{\cod(f)} \circ T(\eta_{\cod(f)}) \circ T(f) = (\mu \circ T \eta)_{\cod(f)} \circ T(f) = (\id_T)_{\cod(f)} \circ T(f) = T(f)$なので$U_T F_T = T$である.

任意の$A \in \mathscr{C}$と$B \in \mathscr{C}_T$に対して$\Mor_{\mathscr{C}_T}(F_T(A),B) \cong \Mor_{\mathscr{C}}(A, T(B)) \cong \Mor_{\mathscr{C}}(A, U_T(B))$となるので\cref{prop: adjunction_iff_hom_natural_iso}により$F_T \dashv U_T$が存在する.
\end{proof}

\begin{definition} \label{def: CategoryOfAdjunctions}
$\mathscr{C}$を圏とし$(T, \eta, \mu)$を$\mathscr{C}$上のモナドとする.\ 圏$\mathbf{Adj}_T$を次のようにして定める.
\begin{itemize}
	\item 対象を随伴$F \dashv U$でモナド$(T, \eta, \mu)$を生成するものとする.
	\item 射$K \colon F \dashv U \to F' \dashv U'$を共変関手$K \colon \cod(F) \to \cod(F')$で$KF = F'$かつ$U'K = U$を満たすものとする.
% https://q.uiver.app/#q=WzAsNixbMSwwLCJcXG1hdGhybXtjb2R9KEYpIl0sWzEsMSwiXFxtYXRocm17Y29kfShGJykiXSxbMCwwLCJcXG1hdGhybXtkb219KEYpIl0sWzQsMCwiXFxtYXRocm17Y29kfShGKSJdLFs0LDEsIlxcbWF0aHJte2NvZH0oRicpIl0sWzMsMCwiXFxtYXRocm17ZG9tfShGKSJdLFsyLDAsIkYiXSxbMCwxLCJLIl0sWzIsMSwiRiciLDJdLFszLDUsIlUiLDJdLFszLDQsIksiXSxbNCw1LCJVJyJdXQ==
\[\begin{tikzcd}
	{\mathrm{dom}(F)} & {\mathrm{cod}(F)} && {\mathrm{dom}(F)} & {\mathrm{cod}(F)} \\
	& {\mathrm{cod}(F')} &&& {\mathrm{cod}(F')}
	\arrow["F", from=1-1, to=1-2]
	\arrow["{F'}"', from=1-1, to=2-2]
	\arrow["K", from=1-2, to=2-2]
	\arrow["U"', from=1-5, to=1-4]
	\arrow["K", from=1-5, to=2-5]
	\arrow["{U'}", from=2-5, to=1-4]
\end{tikzcd}\]
\end{itemize}
この圏をモナド$T$を生成する\term{ずいはんのけん}{随伴の圏}{category of adjunctions}と呼ぶ.
\end{definition}

\begin{proposition}
$\mathscr{C}$を圏とし$(T, \eta, \mu)$を$\mathscr{C}$上のモナドとする.\ クライスリ圏$\mathscr{C}_T$に対して定まる随伴$F_T \dashv U_T$は随伴の圏$\mathbf{Adj}_T$の始対象である.
\end{proposition}

\begin{proof}
$\mathscr{D}$を圏とし$F \colon \mathscr{C} \rightleftarrows \mathscr{D} \colon U$を共変関手とする.\ $F \dashv U \in \mathbf{Adj}_T$とし$\varepsilon \colon FU \Rightarrow \id_{\mathscr{D}}$を$F \dashv U$の余単位,\ $\eta \colon \id_{\mathscr{C}} \to UF$を$F \dashv U$の単位とする.\ 共変関手$K \colon \mathscr{C}_T \to \mathscr{D}$を次のようにして定める.
\begin{itemize}
	\item $A \in \Ob(\mathscr{C})$に対して$K(A) \coloneq F(A)$とする.
	\item $f \in \Mor_{\mathscr{C}}(A, T(B))$に対して$(\varepsilon F)_{B} \circ F(f)$とする.
\end{itemize}
実際,\ 任意の$A \in \Ob(\mathscr{C})$に対して$K(\eta_A) = (\varepsilon F)_{A} \circ F(\eta_A) = (\varepsilon F \circ F \eta)_A = (\id_F)_A = \id_{F(A)} = \id_{K(A)}$である.\ また任意の$f \in \Mor_{\mathscr{C}}(A, T(B))$と$g \in \Mor_{\mathscr{C}}(B, T(C))$に対して$\mu = U \varepsilon F$より$K(g \circ_{\mathscr{C}_T} f) = K(\mu_C \circ T(g) \circ f) = (\varepsilon F)_C \circ F(\mu_C \circ T(g) \circ f) = (\varepsilon F)_C \circ F((U \varepsilon F)_C \circ UF(g) \circ f)$である.\ $\varepsilon$は自然変換なので次の2つの図式が可換になる.
% https://q.uiver.app/#q=WzAsOCxbMCwwLCJGVUZVRihDKSJdLFsxLDAsIkZVRihDKSJdLFswLDEsIkZVRihDKSJdLFsxLDEsIkYoQykiXSxbMywwLCJGVUYoQikiXSxbNCwwLCJGKEIpIl0sWzMsMSwiRlVGVUYoQykiXSxbNCwxLCJGVUYoQykiXSxbMCwxLCJcXHZhcmVwc2lsb25fe0ZVRihDKX0iXSxbMSwzLCJcXHZhcmVwc2lsb25fe0YoQyl9Il0sWzIsMywiXFx2YXJlcHNpbG9uX3tGKEMpfSJdLFswLDIsIkZVKFxcdmFyZXBzaWxvbl97RihDKX0pIiwyXSxbNCw1LCJcXHZhcmVwc2lsb25fe0YoZyl9Il0sWzQsNiwiRlUoRihnKSkiLDJdLFs1LDcsIkYoZykiXSxbNiw3LCJcXHZhcmVwc2lsb25fe0ZVRihDKX0iXV0=
\[\begin{tikzcd}
	{FUFUF(C)} & {FUF(C)} && {FUF(B)} & {F(B)} \\
	{FUF(C)} & {F(C)} && {FUFUF(C)} & {FUF(C)}
	\arrow["{\varepsilon_{FUF(C)}}", from=1-1, to=1-2]
	\arrow["{FU(\varepsilon_{F(C)})}"', from=1-1, to=2-1]
	\arrow["{\varepsilon_{F(C)}}", from=1-2, to=2-2]
	\arrow["{\varepsilon_{F(g)}}", from=1-4, to=1-5]
	\arrow["{FU(F(g))}"', from=1-4, to=2-4]
	\arrow["{F(g)}", from=1-5, to=2-5]
	\arrow["{\varepsilon_{F(C)}}", from=2-1, to=2-2]
	\arrow["{\varepsilon_{FUF(C)}}", from=2-4, to=2-5]
\end{tikzcd}\]
よって$(\varepsilon F)_C \circ F((U \varepsilon F)_C \circ UF(g) \circ f) = \varepsilon_{F(C)} \circ \varepsilon_{FUF(C)} \circ FUF(g) \circ F(f) = \varepsilon_{F(C)} \circ F(g) \circ \varepsilon_{F(g)} \circ F(f) = K(g) \circ K(f)$となるので$K \colon \mathscr{C}_T \to \mathscr{D}$は共変関手になる.

任意の$A \in \Ob(\mathscr{C})$に対して$K	F_T(A) = K(A) = F(A)$であり任意の$f \in \Mor_{\mathscr{C}}(A,B)$に対して$KF_T(f) = \varepsilon_{F(B)} \circ F(\eta_B) \circ F(f) = (\varepsilon F \circ F \eta)_B \circ F(f) = F(f)$なので$KF_T = F$である.\ また任意の$A \in \Ob(\mathscr{C})$に対して$UK(A) = U(F(A)) = T(A) = U_T(A)$であり任意の$f \Mor_\mathscr{C}(A, T(B))$に対して$UK(f) = U((\varepsilon F)_{B} \circ F(f)) = \mu_B \circ T(f) = U_T(f)$なので$UK = U_T$である.\ よって次の2つの図式が可換になる.
% https://q.uiver.app/#q=WzAsNixbMCwwLCJcXG1hdGhzY3J7Q30iXSxbMSwwLCJcXG1hdGhzY3J7Q31fVCJdLFsxLDEsIlxcbWF0aHNjcntEfSJdLFszLDAsIlxcbWF0aHNjcntDfSJdLFs0LDAsIlxcbWF0aHNjcntDfV9UIl0sWzQsMSwiXFxtYXRoc2Nye0R9Il0sWzEsMiwiSyJdLFswLDEsIkZfVCJdLFswLDIsIkYiLDJdLFs0LDUsIksiXSxbNCwzLCJVX1QiLDJdLFs1LDMsIlUiXV0=
\[\begin{tikzcd}
	{\mathscr{C}} & {\mathscr{C}_T} && {\mathscr{C}} & {\mathscr{C}_T} \\
	& {\mathscr{D}} &&& {\mathscr{D}}
	\arrow["{F_T}", from=1-1, to=1-2]
	\arrow["F"', from=1-1, to=2-2]
	\arrow["K", from=1-2, to=2-2]
	\arrow["{U_T}"', from=1-5, to=1-4]
	\arrow["K", from=1-5, to=2-5]
	\arrow["U", from=2-5, to=1-4]
\end{tikzcd}\]
よって$K$は随伴の圏$\mathbf{Adj}_T$の射である.


$K' \colon \mathscr{C}_T \to \mathscr{D}$を$K'F_T = F$かつ$UK' = U_T$を満たす共変関手とすると任意の$A \in \Ob(\mathscr{C})$に対して$K'(A) = K'(F_T(A)) = F = K(F_T(A)) = K(A)$である.\ 任意の$f \in \Mor_{\mathscr{C}}(A, T(B))$に対して$K'(f) \in \Mor_{\mathscr{D}}(F(A), F(B))$であり$UK'(f) \circ \eta_A = U_T(f) \circ \eta_A = \mu_B \circ T(f) \eta_A = \mu_B \circ \eta_{B} \circ f = f$である.\ 同様に$UK(f) \circ \eta_A = f$であり\cref{prop: adjunction_iff_hom_natural_iso}により$f \mapsto U(f) \circ \eta_A$は全単射なので$K'(f) = K(f)$である.\ ゆえに$K' = K$である.
\end{proof}

\begin{proposition}
$\mathscr{C}$を圏とし$(T, \eta, \mu)$を$\mathscr{C}$上のモナドとする.\ アイレンベルグムーア圏$\mathscr{C}^T$に対して定まる随伴$F^T \dashv U^T$は随伴の圏$\mathbf{Adj}_T$の終対象である.
\end{proposition}